\documentclass[12pt]{article}
\usepackage[utf8]{inputenc}
\usepackage[T1]{fontenc}
\usepackage{amsmath, amssymb, amsfonts}
\usepackage{geometry}
\usepackage{hyperref}
\usepackage{enumitem}

\geometry{margin=1in}
\setlength{\parindent}{0pt}
\setlength{\parskip}{0.75em}

\title{Critical Analysis of Schrödinger's Thought Experiment:\\
Correspondence with Information Theory}
\author{Andrei Preece}
\date{Berlin, 2025}

\begin{document}

\maketitle

\begin{abstract}
This article critically re-evaluates Schrödinger's thought experiment through the lens of information theory and quantum mechanics, introducing the concept of entangled flows to analyze the system's participants. We demonstrate that the notion of the cat existing in a ``both alive and dead'' superposition is logically inconsistent and physically untenable. By treating the experiment as an abstract paradox, we argue it lacks direct relevance to physical reality, proposing instead a state of informational uncertainty described as ``neither alive nor dead.''
\end{abstract}

\section{Introduction}

Schrödinger's thought experiment was designed to highlight the absurdity of applying quantum superposition to macroscopic objects. The setup involves a cat placed in a sealed box with a device that releases poison if a uranium atom decays, potentially killing the cat. While the box remains closed, the cat is purportedly in a superposition of ``both alive and dead.'' This article argues that this interpretation violates principles of information theory and macroscopic physics, as quantum and macroscopic states are not directly interchangeable. We introduce the concept of entangled flows to systematically analyze the experiment's participants and resolve its logical contradictions.

\section{Schrödinger's Experiment Setup}

The experiment begins when a cat is placed inside a sealed box containing a device linked to a uranium atom. If the atom decays, the device releases poison, killing the cat. Until the box is opened, the cat's state is described as a superposition of ``alive'' and ``dead.'' This paradox questions the validity of quantum superposition for macroscopic systems, suggesting an absurdity in the theory itself.

\section{Concept of Entangled Flows}

We define ``flow'' as a continuous process that maintains or changes a system's state over time, influenced by internal or external factors. Flows can represent both physical and informational dynamics, preserving a system's integrity until a state transition occurs. In Schrödinger's experiment, three primary flows interact: the Uranium Flow, the Cat Flow, and the Observer Flow, mediated by the poison and the box's lid.

\subsection{Uranium Flow}

The Uranium Flow represents the quantum state of the uranium atom, described by a superposition:
\[
|\psi_{\mathrm{uranium}}\rangle = \alpha |\mathrm{stable}\rangle + \beta |\mathrm{decayed}\rangle, \quad |\alpha|^2 + |\beta|^2 = 1,
\]
where \(\alpha\) and \(\beta\) are complex probability amplitudes. The decay occurs randomly and is independent of the Cat Flow or Observer Flow. Upon decay, the Uranium Flow transitions to a new state, releasing energy without ceasing to exist.

\subsection{Cat Flow}

The Cat Flow encapsulates the macroscopic system of the cat, comprising numerous entangled quantum states of its constituent particles. Unlike the uranium atom, the cat's macroscopic state (``alive'' or ``dead'') cannot exist in superposition due to decoherence. The Cat Flow continues until terminated by the poison, triggered by the uranium atom's decay. The ``alive'' state corresponds to the uranium atom's ``stable'' state, while the ``dead'' state follows the ``decayed'' state.

\subsection{Poison as a Mediating Link}

The poison serves as a critical link between the Uranium Flow and the Cat Flow, entangling their quantum states. It acts as an ``internal observer'' within the box, transmitting information about the uranium atom's decay to the cat's state, thus altering the Cat Flow.

\subsection{Box and Lid}

The box isolates the Uranium-Cat Flow from external influences, while closing the lid initiates the experiment and entangles the Observer Flow with the system. Opening the lid collapses the system's quantum uncertainty, fixing the cat's state and concluding the experiment.

\subsection{Observer Flow}

The Observer Flow encompasses two roles:
\begin{itemize}
\item Internal observer: Conceptually represented by the poison, which links the uranium and cat states.
\item External observer: The entity that opens the box, observes the cat's state, and collapses the system's entanglement.
\end{itemize}

This dual role highlights the observer's simultaneous internal and external relationship to the system.

\section{Decoherence and the Impossibility of Superposition}

The cat's macroscopic state cannot exist in a quantum superposition due to decoherence, where environmental interactions rapidly collapse the system into a definite state. The system's state evolves as:
\[
\rho_{\mathrm{system}} = \sum_i c_i(t)|\phi_i\rangle\langle\phi_i|,
\]
where \(c_i(t)\) represents probability amplitudes that decay exponentially due to environmental interactions:
\[
c_i(t) \propto e^{-\Gamma t},
\]
with \(\Gamma\) as the decoherence rate. This process ensures the Cat Flow settles into either ``alive'' or ``dead'' almost immediately, negating the possibility of a sustained superposition.

\section{Probabilistic Analysis}

The uranium atom's decay follows an exponential probability law:
\[
P_{\mathrm{decay}}(t) = 1 - e^{-\lambda t},
\]
where \(\lambda\) is the decay constant. The cat's state probability, \(P_{\mathrm{cat}}(t)\), depends on the uranium atom's state but is not a quantum superposition. Instead, it reflects informational uncertainty until observation. The system's density matrix, incorporating all flows, is:
\[
\rho_{\mathrm{system}} = \rho_{\mathrm{uranium}} \otimes \rho_{\mathrm{cat}} \otimes \rho_{\mathrm{observer}}.
\]
The poison acts as an interaction operator, transmitting decay information and breaking flow entanglement upon observation.

\section{Reinterpreting the Paradox}

The notion of the cat being ``both alive and dead'' is a misapplication of quantum superposition to macroscopic systems. We propose that the cat's state is better described as ``neither alive nor dead,'' reflecting informational uncertainty rather than a physical superposition. This interpretation aligns with Nāgārjuna's metaphor of a chicken that is neither fully cooked nor raw, illustrating the logical absurdity of opposing states coexisting in macroscopic systems.

\subsection{Logical Advantages}

\begin{itemize}
\item The ``neither alive nor dead'' state avoids contradictions by emphasizing uncertainty over simultaneous opposing states.
\item It aligns with decoherence theory, which precludes macroscopic superpositions.
\item It reframes the experiment as a philosophical inquiry rather than a physical reality.
\end{itemize}

\section{Additional Considerations}

\begin{itemize}
\item Arbitrary Timing: The box's opening is not fixed, requiring analysis of temporal dynamics in the flows.
\item Natural Death: The cat could die naturally before uranium decay, further undermining the superposition concept.
\item Coherence: The ``neither alive nor dead'' model eliminates the need for opposing states, enhancing logical consistency.
\end{itemize}

\section{Conclusion}

Schrödinger's thought experiment does not validate or refute quantum superposition, as it erroneously conflates quantum and macroscopic systems. By applying information theory and the concept of entangled flows, we propose that the cat's state is one of informational uncertainty (``neither alive nor dead'') rather than a physical superposition. This reinterpretation resolves the paradox, aligns with physical reality, and reframes the experiment as a philosophical tool rather than a physical model.

Namaste

\end{document}
